\documentclass[letterpaper]{article}
\usepackage{times}
\usepackage{helvet}
\usepackage{courier}
\usepackage{graphicx}
\usepackage{amsmath}
\usepackage{booktabs}
\usepackage{hyperref}
\usepackage[margin=1in]{geometry}

\title{A Belief-Based Triage Agent for Wearable ECG Alerts:\\
Cost-Asymmetric Decision Making Under Hidden State Uncertainty}

\author{
    Student Preprint\\
    AI Native Shipping Sprint --- Week 1\\
    Belief Engine Course\\
    \texttt{https://github.com/Julie-max/ai-native-sprint}
}

\date{}

\begin{document}

\maketitle

%----------------------------------------------------------
\begin{abstract}
%----------------------------------------------------------
I built an agent that checks a person's ECG from a wearable device
and decides if they need to visit a doctor by reading the ECG ST
segment along with prior history and symptoms of the person wearing
it. The existing approaches are simple static boundaries for alerts
on ECG abnormalities which have high accuracy but miss critical
cases. My agent has a cost-based belief system which is updated
constantly upon new evidence and discovery; it also holds its
decision if there is ambiguity. The thresholds set are: escalate
above 35\% STEMI probability and dismiss below 0.02\% STEMI
probability, derived from a cost asymmetry of \$500,000 for a
missed critical alert versus \$200 for an unnecessary escalation
--- a ratio of 2,500:1.
\end{abstract}

%----------------------------------------------------------
\section{Introduction}
%----------------------------------------------------------
The problem that exists is that every wearable ECG device out there,
and there are many, has simple threshold alarms that treat false
electrode artifacts and critical STEMI alerts as the same case.
Accuracy as a metric hides this failure: a classifier that is 95\%
accurate but always wrong on the critical 5\% is not 95\% useful
--- it is actively dangerous.
My agent works with a cost-based belief system that reasons with
given evidence before making a decision, replacing the dumb
threshold alarm with a three-action belief engine that distinguishes
between hidden states.

%----------------------------------------------------------
\section{Agent Design}
%----------------------------------------------------------
My four inputs for the agent are ST segment ECG morphology, age,
prior medical history, and symptoms experienced. The three hidden
states are True STEMI, Hyperventilation, and Electrode Artifact.
The three actions are: Escalation with an alert to consult a doctor,
Dismiss the alert, or HOLD for further evidence.

The policy is that escalation triggers above 35\% STEMI probability,
dismissal triggers below 0.02\%, and everything in between causes
the agent to HOLD for further evidence. Upon new evidence, the
belief is updated such that the probabilities change and are
re-evaluated against the thresholds.

\subsection*{Threshold Derivation}

Thresholds are derived from expected cost calculations:

\begin{table}[h]
\centering
\begin{tabular}{lll}
\toprule
\textbf{Threshold} & \textbf{Calculation} & \textbf{Result} \\
\midrule
Dismiss (0.02\%) & $0.0002 \times \$500{,}000$ vs $0.9998 \times \$200$ & \$100 $<$ \$199.96 \\
Escalate (35\%) & $0.35 \times \$500{,}000$ vs $0.65 \times \$200$ & \$175{,}000 $\gg$ \$130 \\
\bottomrule
\end{tabular}
\caption{Expected cost derivation for action thresholds}
\end{table}

\subsection*{Cost Model}

\begin{table}[h]
\centering
\begin{tabular}{ll}
\toprule
\textbf{Outcome} & \textbf{Cost} \\
\midrule
Missed true STEMI & \$500,000 \\
Unnecessary escalation & \$200 \\
Cost asymmetry ratio & 2,500:1 \\
\bottomrule
\end{tabular}
\caption{Cost model for the triage agent}
\end{table}

%----------------------------------------------------------
\section{Experiments}
%----------------------------------------------------------
For the experiments, 20 handcrafted test cases have been created,
each with a combination for each hidden state. For example, a case
with 7\% STEMI probability has its belief updated upon feature
change until crossing a threshold for a definite action to be taken
--- starting at HOLD, updating through family history confirmation
(12\%), symptom onset (28\%), and diaphoresis (42\%), at which
point the agent escalates. Most cases land in HOLD as the agent is
conservative by design and appropriately so for this problem
statement.

%----------------------------------------------------------
\section{Discussion}
%----------------------------------------------------------
The discussion on Reddit taught me how the agent I am building is
different from a simple wearable device having a dumb threshold ---
rather, what we have is an agent that updates its belief on new
evidence before acting. One limitation is that the test cases are
simulated and not drawn from actual clinical recordings. Future
directions include using real MIT-BIH Arrhythmia Database data and
defining a maximum HOLD time, since HOLD cannot persist indefinitely
--- if no new evidence arrives within a defined window, the agent
escalates by default, because in cardiac events, indefinite waiting
is itself a dangerous action.

%----------------------------------------------------------
\section*{Acknowledgements}
%----------------------------------------------------------
Reddit community r/BiomedicalEngineers for feedback on cost-ratio
realism and agent framing. Belief Engine Course, Week 1.

%----------------------------------------------------------
\begin{thebibliography}{9}
%----------------------------------------------------------

\bibitem{mitbih}
Moody, G.B. and Mark, R.G. (1992).
\textit{The impact of the MIT-BIH Arrhythmia Database}.
IEEE Engineering in Medicine and Biology Magazine, 20(3), 45--50.
PhysioNet: \url{https://physionet.org/content/mitdb/}

\bibitem{stemi}
Thygesen, K. et al. (2018).
\textit{Fourth Universal Definition of Myocardial Infarction}.
Journal of the American College of Cardiology, 72(18), 2231--2264.

\bibitem{wearable}
Hannun, A.Y. et al. (2019).
\textit{Cardiologist-level arrhythmia detection and classification
in ambulatory electrocardiograms using a deep neural network}.
Nature Medicine, 25, 65--69.

\end{thebibliography}

\end{document}
